% Copyright (c) 2021 Tobias Briones. All rights reserved.
% SPDX-License-Identifier: CC-BY-4.0
%
% This file is part of https://github.com/tobiasbriones/
% cp-unah-mm700-agricultural-soil-sampling-for-data-analysis
%
% This source code is licensed under the Creative Commons Attribution 4.0
% International License found in the LICENSE-CC file in the root directory of
% this source tree or at https://spdx.org/licenses/CC-BY-4.0

\documentclass{report}
\usepackage{framework}

\title{Marco Teórico de Investigación: Modelo de Muestreo Virtual de Suelos Agrícolas en Honduras para Analítica en Base al Muestreo Aleatorio Simple y Estratificado}
\author{Tobias Briones \bigbreak tobias.briones@unah.hn}
\date{Noviembre 2021}

\begin{document}

\makeatletter
    \begin{titlepage}
        \begin{center}
            \includegraphics[width=0.3\linewidth]{ref/logo-unah.png}\\[4ex]
            {\huge \bfseries \@title 
            \vspace{1cm}}\\[2ex]
            {\LARGE \@author}\\[50ex] 
            
            {\large
            Marco Teórico de Investigación Presentado a la\\
            Universidad Nacional Autónoma de Honduras de la Carrera de\\
            Licenciatura en Matemática
            }\\[2ex]
            
            {\large \@date}
        \end{center}
    \end{titlepage}
\makeatother
\thispagestyle{empty}
\newpage

\thispagestyle{empty}
\tableofcontents
\listoffigures
\newpage

\chapter{Marco Teórico}

\import{}{intro}
\import{}{sampling-framework}
\import{}{sampling-size-estimate}
\import{}{srs}
\import{}{sistematic-sampling}
\import{}{stratified-sampling}
\import{}{sampling-patterns}
\import{}{virtual-sampling}
\import{}{agriculture-surveys}
\import{}{hard-sampling}

\chapter{Anexos}

Se ha tomado nota en pequeña parte sobre la elaboración y estructura de \textit{Manual de Elaboración y Presentación de Tesis} \cite{universidad-san-carlos-2016}, \textit{Metodología de la investigación} \cite{collado-2014}, \textit{Introducción a la metodología de la investigación científica} \cite{cabezas-2018} para esta propuesta de investigación.

\section{Glosario}

\begin{itemize}
    \item \textbf{Analítica:} Proceso de detectar, interpretar y comunicar patrones significativos en los datos, así como usar herramientas para que toda la organización pueda realizar cualquier pregunta sobre cualquier información en todos los entornos y dispositivos posibles \cite{oracle-2021}.
    
    \item \textbf{Rendimiento:} Razón con unidad dimensional [MASA][PRODUCTO]/[ÁREA] útil para medir el beneficio producido en una zafra (p. ej. 1TA/HA 1 tonelada de azucar por hectárea).
    
    \item \textbf{Zafra:} Cosecha de la caña dulce (p. ej. zafra 2020-2021).
    
    \item \textbf{Variedad vegetal:} Grupo de plantas definido con mayor precisión, seleccionado dentro de una especie, que presentan una serie de características comunes \cite{union-internacional-para-la-proteccion-de-las-obtenciones-vegatales-upov-2021}.
\end{itemize}

\subsubsection{Términos de Dominio Específico}

\begin{itemize}
    \item \textbf{Usuario:} Empresa agrícola en Honduras.
    
    \item \textbf{Usuario interno:} Empleado de la empresa agrícola.
    
    \item \textbf{Muestreo físico:} Muestra original que el usuario tomó en su suelo agrícola.
    
    \item \textbf{Muestreo virtual:} Resultado o salida del modelo que se ha de desarrollar en este estudio el cual tiene como entrada el muestreo físico.
    
    \item \textbf{Curar:} Transformar los datos originales mediante filtros de forma que se eliminen los todos aquellos datos innecesarios o corruptos para el análisis subyacente.
\end{itemize}

\printbibliography

\end{document}
